\documentclass[10pt]{article}
\usepackage[T1]{fontenc}
\usepackage{lmodern,amsmath,amssymb,amsthm,microtype}
\usepackage[margin=0.85in]{geometry}
\usepackage[colorlinks=true,urlcolor=blue,citecolor=blue]{hyperref}
\newtheorem{theorem}{Theorem}
\newtheorem{lemma}{Lemma}
\title{Exact Counts and Densities for Two Cotangent-Floor Sequences}
\author{Qichao Wang\\{\small Hebei University of Technology}\\[0.3em]
Daoyu Dong\\{\small University of Electronic Science and Technology of China}}
\date{October 26, 2026}
\begin{document}
\maketitle
\begin{abstract}
We prove the two asymptotic conjectures in OEIS A248520 and A248521.
Their counting functions telescope exactly. The increasing-step sequence has
an explicit inverse-cotangent formula; the constant-step sequence has an exact
minimum over at most three consecutive integers. Both asymptotics have bounded
additive error.
\end{abstract}

Put $f(x)=\cot(\pi/(2x))$ for $x\ge1$ and $F(k)=\lfloor f(k)\rfloor$.
Let $a_1(n)$ list the positive integers $k$ with $F(k+1)-F(k)=1$,
and let $a_0(n)$ list those with difference zero. These are A248521 and
A248520 respectively~\cite{one,zero}. Write $\alpha=2/\pi$ and
$\beta=1-\alpha$.

\begin{lemma}\label{lem:step}
Every difference $F(k+1)-F(k)$ is either zero or one. For every $x>1$,
\[
 \alpha(x-1)<f(x)<\alpha x.
\]
\end{lemma}
\begin{proof}
With $u=\pi/(2x)$, differentiation gives
$f'(x)=\alpha(u/\sin u)^2$. On $(0,\pi/2]$, the function
$u/\sin u$ is increasing: its derivative has numerator
$\sin u-u\cos u$, whose derivative is $u\sin u>0$ and whose limit at zero
is zero. Thus $f'$ is decreasing, $f'(x)>\alpha$, and
$f'(x)\le f'(2)=\pi/4<1$ when $x\ge2$.
Consequently $0<f(k+1)-f(k)<1$ for $k\ge2$; also $F(1)=0$ and $F(2)=1$.
These facts prove the difference assertion. Integrating $f'>\alpha$ from
1, where $f(1)=0$, gives the lower bound. The upper bound follows from
$\tan u>u$ for $0<u<\pi/2$.
\end{proof}

\begin{theorem}[Exact counting functions and both conjectures]\label{thm:density}
For $N\ge1$,
\[
 \#\{n:a_1(n)\le N\}=F(N+1),\qquad
 \#\{n:a_0(n)\le N\}=N-F(N+1).
\]
Both sequences are infinite, with
\[
 a_1(n)=\frac\pi2 n+O(1),\qquad
 a_0(n)=\frac\pi{\pi-2}n+O(1).
\]
In particular both asymptotic equivalences conjectured by V. Kotesovec in
the entries are true.
\end{theorem}
\begin{proof}
Lemma~\ref{lem:step} partitions the positive positions into zero and one
steps. Summing the one steps from 1 to $N$ gives $F(N+1)-F(1)=F(N+1)$,
and subtraction gives the zero count. The bounds in the lemma imply
$F(N+1)=\alpha N+O(1)$ and $N-F(N+1)=\beta N+O(1)$.
Both coefficients are positive, proving infinitude. Evaluating the counts
at the $n$-th position of the respective sequence gives
$n=\alpha a_1(n)+O(1)$ or $n=\beta a_0(n)+O(1)$, as required.
\end{proof}

\begin{theorem}[Exact values]\label{thm:exact}
For every $n\ge1$,
\begin{equation}\label{eq:one}
 a_1(n)=\left\lceil\frac\pi{2\arctan(1/n)}\right\rceil-1.
\end{equation}
Put
\[
 L_n=\left\lfloor\frac{n-1}{\beta}\right\rfloor+1,\qquad
 U_n=\left\lfloor\frac n\beta\right\rfloor.
\]
Then
\begin{equation}\label{eq:zero}
 a_0(n)=\min\{k\in\{L_n,\ldots,U_n\}:k-F(k+1)=n\}.
\end{equation}
The candidate interval contains two or three integers, so this is a bounded
explicit formula, not a search over an unbounded sequence prefix.
\end{theorem}
\begin{proof}
The first position $m$ with $F(m)\ge n$ satisfies
$f(m)\ge n$, or $m\ge\pi/(2\arctan(1/n))$.
Since all steps have size at most one, the $n$-th one step occurs at
$m-1$, proving \eqref{eq:one}, including $n=1$.

Write $Z(N)=N-F(N+1)$. The strict bounds on $f$ imply
\[
 \beta N-\alpha<Z(N)<\beta N+1.
\]
At $N=a_0(n)$ the right inequality gives $N>(n-1)/\beta$, so
$N\ge L_n$. Since $U_n>n/\beta-1$, the left inequality gives
$Z(U_n)>n-1$. This integer count is at least $n$, proving $a_0(n)\le U_n$.
The count increases in steps of zero or one, so the first $k$ with count $n$
is exactly the required term. Finally $2<1/\beta<3$, whence
$U_n-L_n+1$ is either 2 or 3.
\end{proof}

\section{Certified evaluation and independent checks}
The verifier encloses $\pi$ using Machin's identity
$\pi=16\arctan(1/5)-4\arctan(1/239)$. It evaluates \eqref{eq:one}
by rational enclosures for $\pi$ and $\arctan(1/n)$, and evaluates $F$
by separate rational sine and cosine
series with alternating remainder bounds. The formula for $a_0$ is then
checked against the literal zero-step definition. All floor and ceiling
operations are certified by matching lower and upper rational bounds;
ambiguous intervals are refined.

For completeness, refinement terminates. The exceptional exact boundaries
$F(1)=0$, $F(2)=1$, and $a_1(1)=1$ are handled symbolically. If
$\cot(\pi/(2k))=m$ were an integer $m\ge2$, then the root of unity
$z=e^{\pi i/k}=(m+i)/(m-i)$ would satisfy
\[
 z+z^{-1}=\frac{2(m^2-1)}{m^2+1}=2-\frac4{m^2+1}.
\]
This rational algebraic integer lies strictly between 1 and 2, a contradiction.
Thus no remaining cotangent floor or inverse-cotangent ceiling lies on an
integer boundary. The linear bounds involving $\beta$ are irrational for
nonzero numerator, since $\pi$ is irrational; the zero numerator is exact.

The supplied official-table prefixes and direct step computations provide
independent finite controls, not the proof of either asymptotic. All arithmetic
uses the Python standard library. The results prove the statements recorded
in the cited entries; they do not assert an exhaustive literature-priority
review or editorial acceptance by OEIS.
\begin{thebibliography}{9}
\bibitem{zero} C. Kimberling, OEIS A248520 (2014),
\url{https://oeis.org/A248520}, accessed October 26, 2026.
\bibitem{one} C. Kimberling, OEIS A248521 (2014),
\url{https://oeis.org/A248521}, accessed October 26, 2026.
\end{thebibliography}
\end{document}
